%!TEX root = ../template.tex
%% Capitulo 3 -- reescrita completa, 2026-09-17.
%% Fontes numericas:
%%   Corpus A (SQ1): LocalResearch/.../bibliometrics/results/corpus_a_registado (331; snapshot 2026-09-13;
%%     re-execucao de controlo corpus_a_2026-09-17 com o mesmo conteudo).
%%   Corpus B (SQ2-SQ3): stage_c/figuras_cap3/tabelas_cap3.py --decisoes
%%     stage_c/verification_v1/verificacao_modelo/decisoes_2026-09-17.json (39 registos; itens pendentes com o valor original).
%%   Fluxo: prisma/PRISMA_FLOW_2026-09-17.md. Correcao B subset A: stage_c/B_SUBSET_A_2026-09-17.md.
%% Versao anterior arquivada em _revisao/2026-09-17-cap3-antes.tex.

\chapter{Revisão de âmbito}
\label{cha:revisao}

Um gestor de energia que abra a literatura à procura de um método para
decidir se a referência de uma unidade ainda serve encontra muitos
indicadores e poucas respostas a essa pergunta. É esse o resultado deste
capítulo. \revisto{Das 331 publicações elegíveis, 162 têm classificação concordante das famílias de modelos; nesse subconjunto, o rácio de intensidade energética é a família mais frequente.} Nas 39 que aplicam
um modelo num processo contínuo dos setores pré-registados, a quantidade
de energia raramente é descrita como medida na fronteira que o modelo
estuda, o domínio de validade quase nunca é declarado e a validação fora
da amostra aparece em quatro artigos. \revisto{Das 130 ocorrências extraídas sobre o que corre mal ao modelo, 89 não nomeiam uma causa física ou operacional.}

O capítulo responde a três perguntas da revisão. A SQ1, sobre a
distribuição dos modelos por setor, tipo de organização e enquadramento
normativo, é tratada sobre o Corpus~A (Secção~\ref{sec:corpus-a}).
A SQ2 e a SQ3, sobre os mecanismos que tornam os modelos inadequados
em processo contínuo e as adaptações propostas, são tratadas sobre o
Corpus~B (Secção~\ref{sec:corpus-b}). A Secção~\ref{sec:gap} reúne o que
falta à literatura e indica onde a dissertação lhe responde. A SQ4 foi
retirada e a razão fica na Secção~\ref{sec:protocolo-sr}.

Convém dizer primeiro o que estes números valem. A extração foi
assistida por modelos de linguagem. No Corpus~A usam-se os valores em
que duas extrações independentes concordam, sem adjudicação humana campo
a campo. No Corpus~B, cada valor tem uma citação ancorada à página,
confirmada por script; o autor analisou e decidiu os casos duvidosos. Nenhum destes procedimentos equivale à
verificação humana integral que o protocolo previa. Por isso, cada
proporção indica o denominador e o nível de verificação, e as
ausências são apresentadas como resultados de pesquisa no texto extraído,
não como propriedades da literatura.

\fecho{Importa precisar o que a concordância entre extrações é e não é. É uma
regra de seleção do que entra na análise, não uma medida de fiabilidade: o
protocolo registado usa-a para encaminhar a atenção humana e adverte
expressamente que não seja tomada como evidência de correção, e é assim que
aqui se usa. Daqui decorre um limite distinto do da cobertura. O subconjunto
concordante não é uma amostra aleatória do corpus: os campos em que os dois
modelos coincidem tendem a ser os que os artigos reportam com mais clareza,
pelo que as distribuições do Capítulo descrevem a parte melhor reportada da
literatura recuperada. Por isso, todas elas trazem o denominador e a categoria
por resolver aparece desenhada em vez de omitida.}

% A abertura do capítulo termina aqui. O espaço livre deve ficar no rodapé,
% sem esticar os intervalos entre os três parágrafos desta página.

\section{Protocolo e desenho da revisão}
\label{sec:protocolo-sr}

\subsection{Porquê uma revisão de âmbito}

A literatura sobre indicadores e referências energéticas mistura estudos
de caso, propostas de método, ferramentas, revisões e documentos de
enquadramento. Não há uma intervenção comum cujo efeito se possa agregar,
e a pergunta não é se um modelo funciona, mas que modelos existem, onde
se aplicam e o que se diz sobre os seus limites. Essa é a situação para a
qual as revisões de âmbito foram propostas: mapear conceitos, tipos de
evidência e lacunas num campo heterogéneo, sem avaliação formal do risco
de viés de cada estudo~\cite{Arksey2005,Levac2010,Munn2018}. O relato segue
a extensão PRISMA-ScR~\cite{Tricco2018}, cuja lista de verificação está no
Apêndice~\ref{app:protocolo-sr}. Seguir a lista organiza o relato; não
certifica a qualidade da revisão.

Três perguntas orientam a análise. A SQ1 pergunta como se distribuem os
modelos de indicador (EnPI) e de referência (EnB) por ano, setor, tipo de
organização e enquadramento normativo. A SQ2 pergunta que mecanismos de
variabilidade operacional tornam esses modelos inadequados em processo
contínuo, e com que efeito sobre o modelo. A SQ3 pergunta que adaptações
são propostas, distinguindo propor de implementar e implementar de avaliar.
Uma quarta pergunta, sobre a aplicabilidade das recomendações numa
refinaria, fazia parte do protocolo e foi retirada a 16 de setembro de 2026
(desvio D8): dependia de uma análise integral às cegas das publicações de
refinação e petroquímica que não foi feita. A Secção~\ref{sec:corpus-b-proximos}
lê criticamente esses estudos, mas não responde à SQ4.

\subsection{Registo, emendas e desvios}

O protocolo foi registado no Open Science Framework a 24 de fevereiro de
2026 (DOI \nolinkurl{10.17605/OSF.IO/TDW6C})~\cite{Ribeiro2026Protocolo},
depois de executadas as pesquisas e antes da recuperação dos textos
integrais e da extração. O registo não protege, portanto, as expressões de
pesquisa contra ajustes feitos depois de ver resultados; protege a
elegibilidade, a extração e a síntese. Seguiram-se uma segunda versão em
abril, a emenda C-2.3 ao protocolo de validação, depositada a 23 de julho,
e a emenda do Stage~C, de 13 de setembro, que fixou o instrumento de
extração do Corpus~B.

A execução afastou-se do registado em nove pontos, listados com data no
Apêndice~\ref{app:protocolo-sr} (desvios D1 a
D9)\matref{; a regra que decide quando um desvio obriga a emenda está na
Equação~\ref{eq:mat-gatilho}}. Quatro pesam nos
resultados deste capítulo e são retomados onde custam alguma coisa: o
Corpus~B foi extraído por uma só família de modelos (D4); a cobertura
humana prevista para esse corpus, 11 publicações às cegas e 31 com
verificação dos campos que sustentam resultados, não foi executada na extensão pretendida (D7);
a SQ4 foi retirada (D8); e o rastreio de referências não foi feito (D9).

\subsection{Dois corpora, duas profundidades}
\label{sec:dois-corpora}

O Corpus~A reúne as publicações que satisfazem os critérios de
elegibilidade em texto integral, qualquer que seja o setor, e sustenta a
SQ1. O Corpus~B é a parte de A que aplica, valida ou demonstra
empiricamente um modelo num processo contínuo de um dos setores fixados
no protocolo antes da seleção: refinação de petróleo, petroquímica,
cimento e cal, pasta e papel, siderurgia e outros metais primários de
produção contínua, química contínua, gases industriais e vidro de fusão
contínua. Uma menção ao setor na introdução ou numa revisão não chega.
B é um subconjunto de A e não uma segunda amostra: cada publicação de B
conta uma vez na bibliometria e é lida com mais profundidade na síntese.

A lista de setores foi fixada antes da seleção por uma razão de desenho.
A presença de incrustação, desativação de catalisadores ou mudanças de
campanha é o fenómeno que a SQ2 quer caracterizar; se servisse para
escolher os estudos, a síntese encontraria o que usou para selecionar.
Pela mesma razão, 13 publicações de processos contínuos fora da lista,
entre elas uma unidade de moagem húmida de milho e um centro de recolha
de petróleo e gás, ficaram como casos exploratórios declarados e não
entram em nenhuma contagem do Corpus~B.

\revisto{A Figura~\ref{fig:desenho-revisao} sintetiza a relação entre os
corpora, as perguntas da revisão e a componente empírica da dissertação.}

\input{5-Figures/tikz/f03a-desenho-revisao}

\section{Pesquisa, seleção e extração}
\label{sec:pesquisa}

\subsection{Fontes e expressões de pesquisa}

A pesquisa correu em fevereiro de 2026 na Scopus e na Web of Science Core
Collection, com acesso pela b-on, e o ficheiro foi exportado para o Rayyan
a 24 de fevereiro. Usaram-se duas expressões. A expressão~A procura o
enquadramento normativo (ISO~50001, ISO~50006 e termos associados); a
expressão~B procura os conceitos de indicador e referência energética em
contexto industrial, sem exigir a norma. Ambas foram limitadas a publicações
posteriores a 1995, em inglês, dos tipos artigo, comunicação em conferência
e revisão. A Scopus devolveu 954 registos (468 pela expressão~A e 486 pela~B)
e a Web of Science 524 (263 e 261), num total de 1\,478. As expressões
literais, a calibração que as fixou e a verificação de sensibilidade com
publicações-âncora estão no Apêndice~\ref{app:protocolo-sr}. Os nomes das
expressões não coincidem com os dos corpora: um registo encontrado pela
expressão~A não fica, por isso, no Corpus~A.

A escolha de vocabulário tem um custo que se consegue medir. O protocolo
previa rastreio de referências para trás e para a frente, com uma regra que obrigava a alargar as expressões se trouxesse
publicações recorrentes que nenhuma apanhava. Esse rastreio não foi feito (desvio D9, 16 de setembro de 2026):
o rendimento é nulo por omissão e não por resultado. Para saber o que a
omissão custa, cruzaram-se as listas de referências dos 42 relatórios então
atribuídos ao Corpus~B com os 925 registos únicos da biblioteca Rayyan,
qualquer que fosse a decisão de triagem. A âncora que o protocolo esperava
recuperar por rastreio~\cite{Karcher2015} já tinha entrado pela pesquisa.
Mas Worrell, citado em 10 das 42 listas, não aparece em nenhum dos 925
registos, e o mesmo com Phylipsen, Bunse, Hasanbeigi e Schulze, citados por três
a sete. É literatura de \emph{benchmarking} e de
gestão de energia industrial que não usa o vocabulário EnPI/EnB em que a
expressão~B se apoia. Com a regra registada, teria obrigado a alargar a pesquisa. O efeito provável recai sobre o enquadramento de que
os estudos incluídos partem, mais do que sobre os modelos cartografados;
não foi testado.

\subsection{Elegibilidade e fluxo}

Títulos e resumos foram avaliados na mesma passagem, no Rayyan, e não em
dois níveis sucessivos como o protocolo previa (desvio D1). Os registos
incluídos ou de decisão incerta passaram à avaliação em texto integral.
O limiar de elegibilidade é metodológico: a publicação tem de apresentar,
propor, avaliar ou criticar um modelo de desempenho energético com detalhe
suficiente para o identificar e caracterizar. Uma referência de passagem a
EnPI ou EnB não chega. A existência de validação não é condição de entrada,
é uma característica a extrair, pelo que uma proposta conceptual ou uma
crítica metodológica podem ser elegíveis.

Dos 1\,478 registos identificados, o Rayyan removeu 553 duplicados. Na
triagem dos 925 restantes excluíram-se 326. Dos 599 enviados para
recuperação, três revelaram-se duplicados tardios e 137 não foram
recuperados, ficando 459 para avaliação em texto integral. Excluíram-se
dez por idioma, um por retratação e 117 por não apresentarem um modelo com
detalhe suficiente. O Corpus~A tem 331 publicações e o Corpus~B 40, das
quais 39 foram cartografadas; uma não pôde ser extraída e a razão fica
documentada no Apêndice~\ref{app:validacao}.

A não recuperação não foi uniforme. Entre os 596 registos que seguiram
para recuperação, obteve-se o texto de 91,6\% dos que tinham DOI e de
30,3\% dos que não tinham; de 85,0\% dos incluídos na triagem e de 65,4\%
dos de decisão incerta. Os textos em falta pesam mais sobre as
comunicações em conferência e a literatura sem identificador, e a sua
elegibilidade é desconhecida: não contam como exclusões por conteúdo.

Uma correção de pertença foi feita a 17 de setembro de 2026, depois do
fecho do corpus. Dois relatórios tinham sido classificados, na adjudicação
às cegas da pertença, como de processo contínuo e, ao mesmo tempo, como não
elegíveis por falta de modelo. Um relatório fora de A não pode estar em B.
Manteve-se a decisão de elegibilidade, tomada antes de qualquer extração, e
os dois saíram do Corpus~B, que passou de 42 para 40. A elegibilidade não
foi reaberta, porque isso significaria decidir a pertença já sabendo o que
cada relatório trazia para a síntese. A regra está hoje imposta no código
que gera os corpora.

A Figura~\ref{fig:prisma} apresenta o percurso completo, da identificação à inclusão.

\begin{figure}[htbp]
  \centering
  \begin{tikzpicture}[x=1cm, y=1cm, font=\scriptsize,
      pb/.style={ntcaixa, text width=6.2cm, align=left, minimum height=10mm},
      px/.style={ntcaixa, text width=5.2cm, align=left, minimum height=10mm,
                 fill=corFundo},
      fase/.style={font=\scriptsize\bfseries, text=corTinta, rotate=90,
                   align=center}]
    \node[pb] (i1) at (4.05,0)
      {Registos identificados nas bases ($n=1\,478$)\\[1pt]
       {\tiny Scopus $n=954$ \quad Web of Science $n=524$}};
    \node[px] (x1) at (11.15,0)
      {Registos removidos antes da triagem\\[1pt]
       {\tiny duplicados $n=553$}};
    \draw[ntseta] (i1) -- (x1);

    \node[pb] (i2) at (4.05,-1.85)
      {Registos triados por título e resumo\\ ($n=925$)};
    \node[px] (x2) at (11.15,-1.85)
      {Registos excluídos ($n=326$)\\[1pt]
       {\tiny fora de âmbito na triagem}};
    \draw[ntsetaf] (i1) -- (i2);
    \draw[ntseta] (i2) -- (x2);

    \node[pb] (i3) at (4.05,-3.55)
      {Relatórios procurados para recuperação\\ ($n=599$)};
    \node[px] (x3) at (11.15,-3.55)
      {Relatórios não recuperados ($n=137$)\\[1pt]
       {\tiny duplicados detetados tardiamente $n=3$}};
    \draw[ntsetaf] (i2) -- (i3);
    \draw[ntseta] (i3) -- (x3);

    \node[pb] (i4) at (4.05,-5.35)
      {Relatórios avaliados para elegibilidade\\ ($n=459$)};
    \node[px, minimum height=14mm] (x4) at (11.15,-5.35)
      {Relatórios excluídos ($n=128$)\\[1pt]
       {\tiny sem modelo de \emph{baseline} $n=117$}\\
       {\tiny idioma não elegível $n=10$ \quad retratação $n=1$}};
    \draw[ntsetaf] (i3) -- (i4);
    \draw[ntseta] (i4) -- (x4);

    \node[pb, draw=corFG, line width=.7pt, fill=corFG!7, minimum height=14mm] (i5)
      at (4.05,-7.35)
      {\textbf{Estudos incluídos --- Corpus~A ($n=331$)}\\[1pt]
       {\tiny Corpus~B, cartografia detalhada: $n=40$ ($39$ cartografados)}\\
       {\tiny casos exploratórios tratados à parte: $n=13$}};
    \draw[ntsetaf] (i4) -- (i5);

    \node[fase] at (0.35,0) {Identificação};
    \node[fase] at (0.35,-3.60) {Triagem};
    \node[fase] at (0.35,-7.35) {Inclusão};
    \draw[corCinzaC, line width=.5pt] (0.80,0.55) -- (0.80,-0.55);
    \draw[corCinzaC, line width=.5pt] (0.80,-1.30) -- (0.80,-5.95);
    \draw[corCinzaC, line width=.5pt] (0.80,-6.65) -- (0.80,-8.05);
  \end{tikzpicture}
  \caption[Fluxograma PRISMA 2020 da seleção de estudos]%
  {Fluxograma PRISMA~2020 da identificação, triagem e inclusão de estudos.
   Não há caixa de rastreio de referências: a pesquisa em bola de neve foi
   registada no protocolo mas não executada (decisão~D9,
   Apêndice~\ref{app:protocolo-sr}). O Corpus~B é um subconjunto do
   Corpus~A, escolhido para cartografia detalhada; os treze casos
   exploratórios são tratados à parte e não entram nas contagens do
   Corpus~A.}
  \label{fig:prisma}
\end{figure}


\subsection{Extração assistida e níveis de verificação}
\label{sec:extracao-verificacao}

A extração foi feita em três camadas, com instrumentos diferentes, e a
confiança que cada número merece depende da camada de onde vem. O
Apêndice~\ref{app:validacao} descreve os instrumentos e os controlos;
aqui interessa o que cada camada permite afirmar.

A pertença aos corpora foi pré-extraída por duas famílias de modelos de
linguagem e as divergências foram adjudicadas pelo autor em formulários às
cegas, sem ver a proposta dos modelos antes de decidir.
Os campos do Corpus~A (tipo de estudo, setor, tipo de organização, norma,
protocolo de medição e verificação, famílias de modelos) foram extraídos
pelas mesmas duas famílias. A análise usa só os valores em que as duas
concordam. As divergências, as extrações únicas e as células vazias ficam
como ``por resolver'', e esse estado aparece nas figuras. Concordância não é
verificação: dois modelos podem ler mal a mesma frase da mesma maneira.

O Corpus~B foi extraído de novo, em setembro, com um instrumento próprio
(esquema v2.3.0-C, congelado a 13 de setembro) e uma só família de modelos
(desvio D4). Cada registo organiza-se em aplicações, modelos e ocorrências.
Um modelo é descrito pela variável dependente, forma, variáveis explicativas,
dados, medição na fronteira, regime de reestimação, domínio de validade e
métricas (campos C1 a C8). Um mecanismo é codificado em dois eixos, a causa
física ou operacional que o artigo nomeia e o efeito sobre o modelo, e
recebe um estatuto de afirmação: \emph{general}, quando é uma afirmação
sobre métodos em princípio (uma definição, uma cautela tirada da
literatura); \emph{applied to the case}, quando é dita sobre o modelo do
próprio estudo; \emph{empirically demonstrated}, quando é mostrada com os
dados do estudo. Uma adaptação é codificada por tipo, estado (proposta,
implementada, avaliada contra uma alternativa), enquadramento normativo e
ligação explícita a um mecanismo. Todos os valores têm uma citação e a
página física de onde vêm, e um script confirma que a citação existe
naquela página.

A verificação humana de cada ocorrência, prevista para este corpus, não foi
executada (desvio D7). Em substituição, a 17 de setembro, as 368 células
com citação de que os resultados deste capítulo dependem foram relidas por
um segundo passe de modelo. O verificador confirmou 225 valores, propôs
correção em 104 e remoção em 21, deu 11 como não sustentados sem valor
alternativo e confirmou outros 7 com baixa confiança. As propostas foram
triadas pelo autor segundo regras fixas: aceitaram-se as que tornam o valor mais
fraco (uma afirmação \emph{general} em vez de aplicada, uma causa não
declarada em vez de nomeada, a remoção de uma limpeza de dados sem
consequência para o modelo) e só se aceitaram as que o tornam mais forte
quando havia evidência quantitativa na página ou citação da norma. Daqui
resultaram 91 correções e 16 remoções. É de notar que o verificador pertence à
mesma família de modelos do extrator, pelo que erros partilhados podem
passar sem sinal. 

\section{Distribuição da literatura (Corpus~A)}
\label{sec:corpus-a}

\subsection{Unidade de análise e tratamento dos dados}

A análise combina bibliometria descritiva (produção, fontes, autoria,
citações) com a distribuição de características extraídas do texto integral,
segundo famílias de análise habituais no mapeamento
científico~\cite{Aria2017Bibliometrix,VanEck2010VOSviewer}. \fecho{A produção e as fontes estão na Secção~\ref{sec:evolucao} e a
geografia na Secção~\ref{sec:geografia}.}

A unidade de contagem é a publicação. Uma publicação pode descrever vários
modelos ou várias instalações, e publicações diferentes podem descrever o
mesmo estudo, pelo que o número de publicações não é o número de casos nem
de validações independentes. Nas variáveis que admitem várias categorias,
cada publicação conta uma vez em cada categoria que lhe foi atribuída, e as
percentagens podem somar mais de 100\%. \novo{Ao longo do capítulo, «campo»
designa uma variável do formulário de extração.}

Os metadados do ficheiro exportado foram completados com o OpenAlex e a
Crossref a 13 de setembro de 2026~\cite{Priem2022OpenAlex,Crossref2026API},
com correspondência por DOI e confronto de títulos. Obtiveram-se metadados
do OpenAlex para 282 das 331 publicações e da Crossref para 278, resumos
para 222 e afiliações para 267. O ficheiro exportado tem autoria, ano e
fonte para todas, palavras-chave para 323, DOI para 299 e nenhum resumo.
A falta de metadados externos não exclui ninguém: publicações sem DOI ou
sem correspondência ficam no conjunto. As contagens de citações das duas
fontes não se fundem.

\alterado{A Figura~\ref{fig:cobertura-campos} fixa o denominador das
distribuições que se seguem: dá, campo a campo, o estado de extração nas
331 publicações do Corpus~A.}

%% 2026-09-21: a figura passou para o fim da 3.3.1, antes da 3.3.2, para
%% imprimir na mesma pagina em que e anunciada e antes das distribuicoes.
%% Com [tbp] e declarada depois do texto da 3.3.2, caia no topo da pagina
%% seguinte, ja depois de a 3.3.3 ter comecado -- e a legenda afirmava uma
%% ordem que a paginacao contradizia. A frase que a legenda tinha sobre isso
%% passou para o corpo, onde a referencia cruzada tem de estar.
\begin{figure}[htbp]
  \centering
  \includegraphics[width=0.93\textwidth]{gerado/cap3-cobertura-campos}
  \caption[Estado de extração campo a campo no Corpus A]{Estado de extração
  de cada campo nas 331 publicações do Corpus~A. Concordância entre modelos
  significa que duas extrações independentes deram o mesmo valor --- não é
  validação humana nem classificação adjudicada. \novo{Divergência é haver dois
  valores diferentes; extração única, só uma das extrações ter produzido valor;
  sem extração, nenhuma o ter feito.} As famílias de modelos, o
  setor e o país são os campos com menos cobertura, e é por isso que os
  resultados dessas secções se leem sobre denominadores menores do que 331.}
  \label{fig:cobertura-campos}
\end{figure}

\subsection{Cobertura dos campos extraídos}

Nos campos lidos do texto integral, a análise usa os valores concordantes
entre as duas extrações (Secção~\ref{sec:extracao-verificacao}). A cobertura
varia muito de campo para campo. Há concordância no tipo de estudo em 258
publicações, no enquadramento normativo em 258, no protocolo de medição e
verificação em 271 e no tipo de organização em 238, \novo{no setor em 201 e
no país em 178}. As famílias de modelos,
o campo mais próximo da SQ1, só têm valor concordante em 162: em 118 as duas
extrações divergem e em 51 só uma produziu valor utilizável. Os 169 casos por
resolver \novo{neste campo} são 51,1\% do corpus.

\novo{A concordância é avaliada sobre a forma normalizada e não sobre a
resposta em bruto: nos campos em texto livre, as respostas passam primeiro
pelo vocabulário controlado que serve a derivação --- os quinze setores de
atividade, os nomes de país da ISO~3166 --- e só depois se comparam. É assim
que \emph{automotive sector} e \emph{automotive manufacturing sector}, ou
\emph{Ecuador} e \emph{coast of Ecuador}, contam como o mesmo valor. Exige-se
igualdade exata da forma normalizada.}

\novo{Nos campos que admitem vários valores, a exigência de igualdade
exata é conservadora noutro sentido: em 91 das 118 divergências do campo das
famílias de modelos as duas extrações partilham pelo menos uma família.}
\corr{Em 64 dessas 91 uma das listas está contida na outra e a diferença é só
de extensão; nas restantes 27 cada extração declara uma família que a outra
não declara, e aí a divergência é de substituição, não de grau. A análise
trata todas como não resolvidas, o que nas 64 subestima a cobertura do campo
em vez de a inflacionar, mas nas outras 27 esconde um desacordo real sobre que
modelo o artigo usa.}

As distribuições seguintes descrevem o subconjunto concordante de cada campo,
com o denominador explícito, e não tratam um valor em falta como ausência da
característica. Um valor concordante entre duas extrações não é uma
classificação adjudicada: a auditoria por
amostragem\matref{~---~cuja definição de erro, intervalos exatos e regra
para o caso de zero erros observados estão nas
Equações~\ref{eq:mat-erro-auditoria} a~\ref{eq:mat-zero-erros}~---}
prevista para estes
campos não foi realizada (Apêndice~\ref{app:protocolo-sr}), e um
contexto pouco representado descreve a cobertura deste corpus, não a
inexistência de métodos nesse contexto.

\fecho{Três observações sobre o instrumento ajudam a ler os números acima. A
primeira: a extração única não está dispersa pelos campos. \corr{Há 81
publicações com pelo menos uma célula nesse estado, mas o grosso} concentra-se
em 51 publicações em que uma das duas superfícies não devolveu praticamente
nada do formulário --- 38 com o campo por preencher em nove dos dez campos e
13 em oito. É por isso que essa categoria vale exatamente 51 em oito dos dez
campos da Figura~\ref{fig:cobertura-campos}\corr{; as restantes 30 têm uma só
célula em extração única, quase todas no país, e é isso que levanta esse campo
aos 68 da figura}. O que ela mede é uma falha de
instrumento e não uma característica das publicações.}

\fecho{A segunda: a análise de palavras-chave não produz temas legíveis. Os
grupos obtidos por otimização de
modularidade\matref{~(Equações~\ref{eq:mat-modularidade}
e~\ref{eq:mat-tfidf})} reúnem, em dois casos, o
vocabulário com que o corpus foi recuperado --- eficiência energética,
indicador de desempenho energético, sistema de gestão de energia, ISO~50001
--- e encontrá-lo agrupado é consequência da expressão de pesquisa e não um
resultado sobre a literatura; num terceiro caso reúnem variantes morfológicas
do mesmo termo. Os rótulos automáticos também não são estáveis: numa
reexecução de controlo, um dos grupos trocou de
rótulo por um empate na quinta palavra-chave. Por isso não se nomeiam aqui
grupos temáticos.}

\fecho{A terceira: as redes bibliográficas cobrem pouco. O
acoplamento\matref{~(Equação~\ref{eq:mat-acoplamento})}, que
liga publicações com pelo menos duas referências em comum, resolve-se em 65
das 331 publicações, 19,6\%; as restantes não partilham referências
suficientes com nenhuma outra depois da resolução de metadados.}

\subsection{Evolução temporal e fontes}
\label{sec:evolucao}

As 331 publicações vão de 1997 a 2026. São de 2015 ou posteriores 281
(84,9\%) e de 2020 ou posteriores 168 (50,8\%); o máximo anual é 35, em 2020,
seguido de 34 em 2025. O único registo de 2026 pertence a um ano coberto
apenas até fevereiro. \alterado{A Figura~\ref{fig:evolucao-anual} mostra a
distribuição anual por tipo de documento e as revistas mais frequentes.}


\begin{figure}[htbp]
  \centering
  \includegraphics[width=\textwidth]{gerado/cap3-evolucao-anual}
  \caption[Evolução anual e revistas do Corpus A]{Publicações do Corpus~A,
  1997--2026. (A) Distribuição anual por tipo de documento; a barra cinzenta
  indica 2026, coberto apenas até fevereiro. As linhas tracejadas assinalam
  as publicações da ISO~50001:2011 e da ISO~50006:2014 como marcos temporais,
  sem estabelecer uma relação causal com a evolução observada. (B) Revistas
  com pelo menos cinco artigos entre os 214 artigos de revista, após
  harmonização das variantes verificadas por ISSN.}
  \label{fig:evolucao-anual}
\end{figure}

\newpage

São artigos de revista 214 publicações (64,7\%) e comunicações em
conferência 117 (35,3\%). O peso das conferências descreve o conjunto e não
diz nada sobre a qualidade de cada estudo. \alterado{Entre os 214 artigos de
revista, contam-se 112 revistas depois de harmonizar as variantes verificadas
por ISSN. Oitenta e três revistas surgem uma única vez; as nove com pelo menos
cinco artigos reúnem 83 (38,8\%). A \emph{Energies} tem 24 artigos, o
\emph{Journal of Cleaner Production} 12 e a \emph{Energy} e a
\emph{Sustainability} 9 cada.}

\subsection{\novo{Geografia dos casos}}
\label{sec:geografia}

\fecho{O país registado é o do caso estudado e não o da afiliação dos autores,
que são coisas distintas e frequentemente confundidas na bibliometria.} \novo{O
país do caso está atribuído em 154 das 331 publicações (46,5\%), em
161 atribuições por 50 países, já que uma publicação pode estudar casos em mais
do que um. A Figura~\ref{fig:mapa-casos} mostra-as. A Itália reúne 21
atribuições e a Colômbia 15; nenhum outro país passa de sete. O que a figura
descreve é a cobertura deste corpus e não a distribuição da literatura: as 177
publicações sem país atribuído não estão no mapa, e uma região sem pontos é
ausência de localização atribuída, não ausência de estudos. Catorze dos 50
países têm pelo menos um caso no Corpus~B. As duas maiores contagens
concentram-se em poucos grupos: um autor assina dez das 21 publicações com
caso em Itália e outro cinco das quinze colombianas. A figura conta
publicações, não grupos de investigação independentes.}

\begin{figure}[htbp]
  \centering
  \includegraphics[width=\textwidth]{gerado/cap3-mapa-casos}
  \caption[País atribuído aos casos do Corpus A]{\novo{País atribuído aos casos
  do Corpus~A. O tamanho do marcador é o número de atribuições ao país e o tom
  distingue os países com pelo menos um caso no Corpus~B. O marcador está no
  ponto de etiqueta do país, não na instalação estudada. Uma região sem
  marcadores significa que nenhuma publicação com localização atribuída a
  refere, e não que não exista literatura sobre ela: 177 das 331 publicações
  não têm país do caso atribuído e não estão representadas. As contagens são de
  publicações e não de grupos independentes: nos dois países com mais
  atribuições, um só autor assina dez das 21 publicações italianas e cinco das
  quinze colombianas.}}
  \label{fig:mapa-casos}
\end{figure}

\subsection{Estudos e contextos}

Das 258 publicações com tipo de estudo concordante, 153 são um caso aplicado
único e 53 estudos com vários casos; há 18 revisões, 16 descrições de
ferramentas, 14 trabalhos metodológicos ou teóricos e quatro modelos de
maturidade. \corr{A aplicação em mais do que uma instalação
é um campo distinto, com valor concordante em 268 publicações; em 17 dessas a
pergunta não se aplica e, das 251 restantes, só 27 aplicam o modelo em mais do
que uma instalação.} A literatura é, em larga medida, uma
coleção de casos isolados, e o trabalho sobre o método propriamente dito é
uma fração pequena.

No tipo de organização, as 238 classificações concordantes repartem-se por
manufatura discreta (71), processo contínuo (58), edifícios (57), contextos
mistos (21), serviços (17) e não especificado (14). A etiqueta ``processo
contínuo'' é mais larga do que a pertença ao Corpus~B, que exige também um
dos setores pré-registados e uma aplicação empírica. O setor e o país dos
casos, em texto livre, não permitem ainda uma distribuição fiável: \novo{130
e 153} publicações, respetivamente, ficam por classificar depois da
normalização. \novo{A distribuição por país está na
Secção~\ref{sec:geografia}.} \corr{Os dois números do país não são o mesmo: as
153 são as que não têm valor concordante, ao passo que as 177 sem país
atribuído dessa secção incluem ainda 24 em que as duas extrações concordam em
não haver localização a atribuir.} \novo{Nenhuma das duas sustenta uma conclusão sobre
concentração setorial ou geográfica da literatura, porque em ambas a categoria
mais frequente é a ausência de classificação.}

%% 2026-09-21: resolve a caixa C3-V05. A parte geografica esta na 3.3.4.
\novo{A Figura~\ref{fig:estudos-contextos} dá os dois campos publicação a
publicação, com a categoria por resolver à vista.}

\begin{figure}[htbp]
  \centering
  \includegraphics[width=\textwidth]{gerado/cap3-estudos-contextos}
  \caption[Tipo de estudo e tipo de organização no Corpus A]{\novo{Tipo de
  estudo e tipo de organização nas 331 publicações do Corpus~A. Cada quadrado é
  uma publicação, pelo que os dois painéis são diretamente comparáveis.}
  \fecho{As categorias estão agrupadas quatro a quatro; as contagens das seis
  categorias de cada campo estão no texto acima.} \novo{Os quadrados vazios
  são as publicações cujo campo ficou por resolver ou sem extração: 73 no tipo
  de estudo e 93 no tipo de organização. O bloco dos casos aplicados únicos
  ocupa quase metade do primeiro painel.}}
  \label{fig:estudos-contextos}
\end{figure}

\subsection{Famílias de modelos e enquadramento normativo}
\label{sec:familias-normas}

\fecho{A Tabela~\ref{tab:a-familias} dá as famílias declaradas pelas 162
publicações com o campo concordante e, ao lado, a combinação entre as duas que
esta dissertação distingue.} O rácio simples de intensidade energética é a
família mais frequente: um rácio ou uma média de consumos específicos aparece
em 90 das 162 publicações e uma regressão em 80. \fecho{São 82 as que não
ajustam nenhuma regressão, e 59 dessas usam um rácio ou uma média.} As
categorias não são exclusivas, e a presença de duas famílias no mesmo artigo
não quer dizer que tenham sido comparadas nem combinadas.

A diferença entre as duas primeiras famílias é o que interessa a esta
dissertação. \revisto{Um rácio divide o consumo pela produção. Usá-lo como normalizador invariável à carga pressupõe proporcionalidade entre as duas grandezas, sem termo constante relevante; a mera divisão não testa esse pressuposto.} Uma regressão
estima esse termo e o declive a partir dos dados, e com isso pode ser
avaliada. A categoria de deteção de mudança e controlo estatístico junta
procedimentos com funções distintas; a sua presença não mostra que uma
referência seja atualizada nem que a não estacionariedade seja tratada.
\fecho{São as três primeiras famílias da tabela --- rácio, regressão linear e
regressão múltipla --- que enquadram o caso real; o
Capítulo~\ref{cha:diagnostico} compara as regressões simples e múltipla, além
dos comparadores triviais.}

%% 2026-09-21: a Figura 3.7 era um ranking de dez categorias que o paragrafo
%% acima ja dava por extenso, e nao mostrava a distincao que a subseccao
%% argumenta. Passa a tabela: os mesmos numeros lem-se melhor em coluna, e o
%% bloco da direita da a particao exclusiva, que a figura nao dava. A
%% v06_modelos() fica em _gerador/figuras.py, sem uso.
{\footnotesize
\ntzebra
\begin{xltabular}{\textwidth}{@{}>{\raggedright\arraybackslash}X r@{\hspace{2.2em}}>{\raggedright\arraybackslash}X r@{}}
\caption[Famílias de modelos no Corpus A]{\fecho{Famílias de modelos de
desempenho energético declaradas pelas 162 publicações do Corpus~A com o campo
concordante; 169 ficaram por resolver ou sem extração. As famílias não são
exclusivas e a coluna da esquerda não soma 162: 84 publicações declaram uma só
família, 62 declaram duas e 16 declaram três ou mais. O bloco da direita é
exclusivo e soma 162; a última linha são as publicações que não declaram
rácio, média nem regressão.}}\label{tab:a-familias}\\
\hiderowcolors
\toprule
\textbf{Família} & \textbf{Publ.} & \textbf{Combinação (exclusiva)} & \textbf{Publ.} \\
\midrule
\endfirsthead
\toprule
\textbf{Família} & \textbf{Publ.} & \textbf{Combinação (exclusiva)} & \textbf{Publ.} \\
\midrule
\endhead
\midrule
\rowcolor{white}\multicolumn{4}{r}{\emph{continua na página seguinte}}\\
\endfoot
\bottomrule
\endlastfoot
\showrowcolors
    Rácio de intensidade & 74 & Rácio ou média, sem regressão & 59 \\
    Regressão linear & 48 & Regressão, sem rácio nem média & 49 \\
    Regressão múltipla & 41 & Rácio ou média e regressão & 31 \\
    Modelo físico & 21 & Só outras famílias & 23 \\
    \cmidrule(l){3-4}
    Ponto de mudança ou CEP & 20 & \emph{Total} & \emph{162} \\
    Aprendizagem automática & 19 & & \\
    Média do consumo específico & 17 & & \\
    Outra & 15 & & \\
    Índice composto & 2 & & \\
    Integração de processo & 2 & & \\
\end{xltabular}
}

%% 2026-09-21: fecha a caixa C3-V07. A normalizacao do setor foi revista e
%% ainda assim a matriz nao sustenta leitura por setor; dito no texto.
\fecho{O cruzamento entre setor e família de modelos não sustenta leitura por
setor. Mesmo depois de revista a normalização, 58 das \corr{120} células do
cruzamento \corr{(doze setores por dez famílias)} têm pelo menos uma
publicação, o máximo é dezasseis e a mediana das células não vazias é dois; a
única linha com substância é a dos edifícios e serviços públicos.}
\corr{Nas etiquetas setoriais que correspondem aos setores de processo
contínuo pré-registados, as contagens vão de uma publicação na refinação de
petróleo e duas na pasta e papel a oito no cimento e cal, oito no vidro e
cerâmica, onze na petroquímica, treze na química e farmacêutica e dezasseis
nos metais; repartidas por dez famílias, nenhuma dessas linhas chega para uma
distribuição.} \fecho{A matriz não se reproduz.}

O enquadramento normativo tem classificação concordante em 258 publicações.
A ISO~50001 aparece em 185 (71,7\%) e a ISO~50006, a norma que trata
especificamente de indicadores e referências, em 48 (18,6\%); seis referem a
ISO~50002, 29 outros referenciais nomeados e 66 nenhum\novo{; as categorias
não são exclusivas, e 65 publicações nomeiam mais do que um referencial, pelo
que os números somam mais do que 258}. Nos protocolos de
medição e verificação, com 271 classificações concordantes, 249 não
identificam nenhum; aparecem o IPMVP em 12, a ASHRAE Guideline~14 em nove e o SEP~M\&V em sete\novo{, o único dos três que fixa critérios
quantitativos de validade do modelo dirigidos a instalações industriais
(Secção~\ref{sec:assimetria-comparacao})}. Esta codificação descreve o que o artigo invoca, não a
certificação da organização nem a conformidade do modelo. E a ausência de
protocolo não quer dizer ausência de medição: quer dizer que o artigo não
se apoia num protocolo reconhecido para a dizer.

%% 2026-09-21: resolve a caixa C3-V08. A caixa pedia as duas marginais; o
%% cruzamento responde a pergunta que as marginais nao respondem, e vai por
%% inteiro para a Tabela B.4, no Apendice B, onde as celulas pequenas nao
%% custam nada. A figura que o dava em heatmap foi retirada: tinha 24
%% celulas, seis a zero e doze com menos de cinco publicacoes, e a cor nao
%% codificava nada que os numeros impressos nas celulas nao dissessem. Um
%% prototipo alternativo -- as duas taxas com intervalo de Wilson -- foi
%% tambem posto de lado: o elemento dominante da imagem eram os intervalos
%% sobrepostos, e isso da peso visual a um nao-resultado. O par de taxas fica
%% em texto, com o travao. A v08_norma() fica em _gerador/figuras.py, sem uso.
\novo{O cruzamento entre os dois campos responde à pergunta que as duas
distribuições separadas não respondem.} \corr{Invocam alguma norma da família
ISO~50001 --- a 50001, a 50002 ou a 50006 --- 191 publicações, mais do que as
185 que nomeiam a ISO~50001 propriamente dita. Dessas 191, 187 têm o campo do
protocolo resolvido; nessas 187, 173 não invocam protocolo algum de medição e
verificação e 14 invocam. As quatro restantes ficam por resolver e não contam
para nenhum dos lados.} \fecho{A proporção que invoca um protocolo é de 7,5\%
--- 14 em 187, intervalo de
Wilson a 95\% de 4,5\% a 12,2\% --- e de 7,8\% entre as 66 que declaram não
invocar norma nenhuma, 5 em 64, de 3,4\% a 17,0\%. Invocar a norma não traz
consigo um procedimento de verificação. Com catorze e cinco casos,
isto descreve o corpus e não testa uma diferença. O cruzamento completo está
na Tabela~\ref{tab:apB-norma-protocolo}.}

%% 2026-09-21: a antiga 3.3.7 ("Estrutura tematica, colaboracao e
%% referencias") foi dissolvida. Descrevia o metodo de redes que o capitulo
%% nao mostra e cujos resultados nao alimentam a SQ1 nem a sintese do Corpus
%% B. Os dois achados que valiam -- a circularidade dos grupos tematicos e a
%% cobertura de 19,6% do acoplamento -- sao sobre o instrumento e passaram
%% para o fim da 3.3.2, junto dos outros limites do instrumento. Fecha
%% tambem as caixas C3-V09 e C3-V10.

\subsection{O que a SQ1 permite dizer}
\label{sec:sq1-sintese}

O protocolo registou três expectativas para a SQ1 antes da extração, e vale
a pena confrontá-las com o que se encontrou, porque duas só se confirmam em
parte.

A primeira dizia que o indicador dominante é um rácio sem ajuste e que as
referências ajustadas são minoria. A primeira metade confirma-se: o rácio é
a família mais frequente. A segunda não: no subconjunto concordante, uma
regressão aparece em 80 de 162 publicações, praticamente tantas como os
rácios. \corr{Quantas dependem apenas de rácios ou médias depende de como se
define a dependência: 59 não ajustam regressão nenhuma e, dessas, 47 não
declaram mais nenhuma família (Tabela~\ref{tab:a-familias}). Em qualquer das
leituras é menos de metade das 162.} \fecho{É o único ponto em
que os dados contrariam uma expectativa registada antes da extração}\corr{; mas
é também o campo com menos cobertura do corpus, pelo que a contradição assenta
em evidência mais fina do que as confirmações, e não mais robusta}. Com metade das
publicações por classificar neste campo, não se pode dizer mais do que isto.

A segunda dizia que a ISO~50001 é quase universal como referência e que a
ISO~50006 e os protocolos de medição e verificação não são. Confirma-se, com
a ressalva de que ``quase universal'' é forte para 185 em 258: é a norma de
gestão que dá nome ao campo, e a norma que diz como construir e manter um
indicador aparece num artigo em cada cinco. \fecho{O contraste mais forte do
capítulo está aqui: \corr{das 187 publicações que invocam a família ISO~50001
e têm o campo do protocolo resolvido, 173 --- mais de nove em cada dez ---}
não invocam protocolo algum de medição e verificação.
A literatura que adota o vocabulário da norma de gestão não adota, na mesma
proporção, nenhum quadro reconhecido para verificar o que mede. É, dentro
deste corpus, a contrapartida empírica do que a Secção~\ref{sec:assimetria-comparacao}
estabelece a partir do texto das próprias normas.}

A terceira dizia que a evidência é feita de casos únicos, em escala de
instalação, geograficamente concentrada, e que o trabalho metodológico é uma
fração pequena. Os casos únicos (153 em 258) e o peso reduzido do trabalho
metodológico (14 em 258) confirmam-se. A escala de aplicação não foi
extraída no Corpus~A\novo{. Quanto à concentração geográfica, a Secção~\ref{sec:geografia}
mostra que as atribuições disponíveis se repartem por 50 países e que as duas
maiores se concentram em poucos grupos de autores; com 177 publicações sem
país atribuído, isso descreve a cobertura do corpus e não permite confirmar
nem infirmar a expectativa}. \corr{O que essa concentração mostra não é, em
rigor, geográfico: dez das 21 publicações com caso em Itália e cinco das
quinze colombianas saem de um autor cada, pelo que o número de casos
independentes é menor do que o número de publicações. É a expectativa dos
casos únicos a confirmar-se por outra via.}

\corr{Das quatro dimensões que a SQ1 enuncia, a distribuição por ano está na
Secção~\ref{sec:evolucao} e não se repete aqui; a do setor fica por responder.
O campo tem valor concordante em 201 das 331 publicações, mas entre as
classificadas a categoria mais frequente é o balde residual, e o cruzamento
com as famílias distribui 58 células não vazias com mediana dois
(Secção~\ref{sec:familias-normas}). Não é que a literatura não varie com o
setor: é que este corpus, extraído deste modo, não permite saber se varia. A
dimensão setorial da SQ1 não é respondível aqui, e isso é um resultado sobre o
instrumento --- a distribuição por setor depende de um campo em texto livre
que a extração assistida não normaliza de forma fiável --- e não uma omissão
da revisão. É também a razão por que o Capítulo~\ref{cha:diagnostico} não pode
apoiar-se em precedentes setoriais e tem de estabelecer os seus comparadores a
partir do caso.}

Para a pergunta que abre o capítulo, \novo{esta cartografia da literatura} diz
sobretudo onde procurar.
A literatura que usa o vocabulário da norma é grande e dispersa, faz casos
e raramente pergunta como se valida uma referência. As 39 publicações do
Corpus~B permitem ir mais fundo do que etiquetas: que variáveis normalizam o
consumo, se a energia é medida onde o modelo diz, se o modelo é refeito e
testado fora dos dados que o ajustaram.

%% 2026-09-21: o titulo era "Sintese do Corpus B", que nao e paralelo ao
%% da 3.3 e roubava a palavra "sintese" a 3.5, que vinha logo a seguir
%% com o mesmo termo. Esta seccao nao sintetiza, examina; e o que examina
%% sao as condicoes em que um modelo pode servir de referencia, que e a
%% frase com que a 3.5 abre.
\section{\ajuste{Condições de uso dos modelos} (Corpus~B)}
\label{sec:corpus-b}

Esta secção lê as 39 publicações cartografadas do Corpus~B. \ajuste{Todos os
números vêm da extração própria deste corpus e da verificação que se lhe
seguiu, ambas descritas na Secção~\ref{sec:extracao-verificacao}}; os denominadores mudam consoante a
unidade de contagem, que é sempre dita: artigos (39), aplicações de um
modelo a um contexto (46), modelos (140) ou ocorrências de um mecanismo ou de
uma adaptação. Nas tabelas, ``n.r.'' é não reportado no texto extraído, e
não reportado não é o mesmo que ausente.

\subsection{Composição}

A composição do corpus condiciona tudo o que se segue
(Tabela~\ref{tab:b-composicao}). Cimento e siderurgia dão 20 das 46
aplicações; refinação e petroquímica dão 12, e há mais duas em artigos com
vários setores. Treze artigos têm pelo menos uma aplicação empírica numa
refinaria ou numa instalação petroquímica, mas nem todos modelam uma
unidade de processo: um é uma central de cogeração, outro um índice de
maturidade do sistema de gestão. A escala dominante é a instalação ou o
conjunto de instalações (31 aplicações); só oito modelam uma unidade e
duas um equipamento. As publicações vão de 2011 a 2025. A maior parte da
evidência sobre mecanismos está, portanto, fora da refinação, e um número
tirado só dos estudos de refinação é um número sobre esses estudos, não
sobre a literatura.

{\footnotesize
\ntzebra
\begin{xltabular}{\textwidth}{@{}>{\raggedright\arraybackslash}X r r r@{}}
\caption[Composição do Corpus B por setor]{Composição do Corpus~B por
setor registado. Um artigo pode ter aplicações em mais do que um setor.
A última coluna conta as aplicações com modelo aplicado empiricamente numa
refinaria ou instalação petroquímica.}
\label{tab:b-composicao}\\
\hiderowcolors
\toprule
\textbf{Setor} & \textbf{Aplicações} & \textbf{Artigos} & \textbf{Refinaria ou petroquímica} \\
\midrule
\endfirsthead
\toprule
\textbf{Setor} & \textbf{Aplicações} & \textbf{Artigos} & \textbf{Refinaria ou petroquímica} \\
\midrule
\endhead
\midrule
\rowcolor{white}\multicolumn{4}{r}{\emph{continua na página seguinte}}\\
\endfoot
\bottomrule
\endlastfoot
\showrowcolors
    Cimento e cal & 10 & 9 & 0 \\
    Siderurgia e metais primários & 10 & 7 & 0 \\
    Petroquímica & 7 & 7 & 7 \\
    Vários setores & 6 & 6 & 2 \\
    Refinação de petróleo & 5 & 5 & 5 \\
    Pasta e papel & 3 & 3 & 0 \\
    Vidro & 2 & 2 & 0 \\
    Química contínua & 2 & 2 & 0 \\
    Gases industriais & 1 & 1 & 0 \\
    \midrule
    Total & 46 & 39 & 14 \\
\end{xltabular}
}

\subsection{Os modelos e o que os alimenta}

Os 39 artigos descrevem 140 modelos. Rácios de intensidade (51 modelos em
19 artigos) e regressões lineares (46 em 19) são dois terços do total; o
resto reparte-se por médias de consumo específico, regressões não lineares,
índices compostos, balanços energéticos, cartas de controlo e casos únicos
de redes neuronais, árvores de decisão, fronteira estocástica e análise
envoltória de dados (Tabela~\ref{tab:b-modelos}). A variável dependente é o
consumo específico em 61 modelos e o consumo absoluto em 48. Onde há
estimação, é quase sempre por mínimos quadrados, e 44 modelos não dizem como
foram estimados.

A carga domina as variáveis explicativas. A quantidade produzida ou
processada aparece em 35 dos 39 artigos, como variável explicativa ou como
denominador, e em 78 dos 140 modelos é a única variável que normaliza o
consumo. Modo de operação, condições ambientais, gama de produtos e
qualidade da alimentação aparecem, cada uma, em nove artigos ou menos; a
qualidade da alimentação, que numa refinaria muda com o crude, em quatro.
Quando há seleção de variáveis, é por juízo do autor (19 modelos) ou não
existe (30); testes de significância e seleção passo a passo contam-se em
sete modelos.

Os dados vêm de auditorias (49 modelos em 14 artigos) quase tanto como de
séries contínuas (45 em 16), e a frequência de modelação mais comum é um
ponto por instalação ou por ano. Nos modelos que o declaram, a série de
ajuste tem uma mediana de 14 meses (57 modelos) e 29 observações (52).
É pouco para distinguir uma mudança da relação de um ano atípico.

{\footnotesize
\ntzebra
\begin{xltabular}{\textwidth}{@{}>{\raggedright\arraybackslash}X r r@{\hspace{2.2em}}>{\raggedright\arraybackslash}X r r@{}}
\caption[Modelos e variáveis explicativas no Corpus B]{Famílias de modelos
e tipos de variável explicativa no Corpus~B. Um modelo pode ter mais do que
uma família. Variáveis: ocorrências com papel de preditor, denominador,
segmentação ou condição de referência.}
\label{tab:b-modelos}\\
\hiderowcolors
\toprule
\textbf{Família} & \textbf{Mod.} & \textbf{Art.} & \textbf{Tipo de variável} & \textbf{Ocorr.} & \textbf{Art.} \\
\midrule
\endfirsthead
\toprule
\textbf{Família} & \textbf{Mod.} & \textbf{Art.} & \textbf{Tipo de variável} & \textbf{Ocorr.} & \textbf{Art.} \\
\midrule
\endhead
\midrule
\rowcolor{white}\multicolumn{6}{r}{\emph{continua na página seguinte}}\\
\endfoot
\bottomrule
\endlastfoot
\showrowcolors
    Rácio de intensidade & 51 & 19 & Carga ou produção & 120 & 35 \\
    Regressão linear & 46 & 19 & Modo de operação & 60 & 9 \\
    Média de consumo específico & 10 & 4 & Outra & 54 & 18 \\
    Regressão não linear & 9 & 5 & Condições ambientais & 27 & 7 \\
    Índice composto & 5 & 5 & Gama de produtos & 18 & 7 \\
    Outra & 5 & 3 & Configuração do equipamento & 11 & 5 \\
    Balanço energético & 3 & 2 & Taxa de utilização & 9 & 6 \\
    Carta de controlo & 3 & 3 & Qualidade da alimentação & 4 & 4 \\
    Rede neuronal & 3 & 1 & Calendário & 1 & 1 \\
    Envoltória de dados (DEA) & 3 & 1 & & & \\
    Outra aprendizagem automática & 2 & 1 & & & \\
    Árvores de decisão & 1 & 1 & & & \\
    Híbrido físico-estatístico & 1 & 1 & & & \\
    Fronteira estocástica & 1 & 1 & & & \\
\end{xltabular}
}

\subsection{Fronteira, validade temporal e validação}
\label{sec:b-fronteira-validade}

A Figura~\ref{fig:mapa-evidencia} mostra, artigo a artigo, sete perguntas que
decidem se um modelo pode servir de referência: se a energia é medida na
fronteira do modelo, se o modelo é refeito, se se diz onde vale, se se
reporta o ajuste, se se testa fora dos dados de ajuste, se as adaptações
foram implementadas ou avaliadas e se dados ou código estão disponíveis.
Cada artigo aparece com o nível mais forte entre os seus modelos, o que
favorece os artigos com vários modelos.

%% 2026-09-21: dois paineis lado a lado. A versao de uma coluna media
%% 6,3 x 10,6 pol; com height=0.86\textheight e keepaspectratio o limite era
%% a altura, a figura ficava com 60% da largura da mancha e o texto interno
%% caia para ~4 pt -- era o que a caixa F05 perguntava. Sai o height, sai o
%% [p], e a legenda passa de dez linhas para cinco: era a legenda que
%% empurrava o folio para cima da ultima linha.
\begin{figure}[htbp]
  \centering
  \includegraphics[width=\textwidth]{cap3/mapa-evidencia-corpus-b}
  \caption[Mapa de evidência do Corpus B]{\ajuste{Mapa de evidência do
  Corpus~B: 39 artigos por sete condições de uso de uma referência, cada
  célula no nível mais forte declarado entre os modelos do artigo, o que
  favorece os artigos com vários modelos. Cinzento é não reportado no texto
  extraído e o tracejado é não aplicável; as definições de cada nível estão no
  Apêndice~\ref{app:validacao}. $\blacklozenge$ aplicação empírica em
  refinaria ou petroquímica.}}
  \label{fig:mapa-evidencia}
\end{figure}

%% 2026-09-21: fecha a caixa F05. As etiquetas nao se liam porque a figura
%% era escalada pela altura; resolvido com os dois paineis. A figura foi
%% regerada COM --decisoes: a que estava impressa tinha sido gerada sem elas
%% e dava 11/39 na reestimacao e 10/39 no dominio, contra os oito e os sete
%% artigos que o texto desta seccao afirma. Versao antiga arquivada em
%% _revisao/2026-09-21r-mapa-evidencia-sem-decisoes.pdf.

A medição na fronteira é a primeira falha. As 140 descrições de modelo têm
163 entradas de medição, uma por vetor energético. Só 35 descrevem uma
correspondência exata entre o que se mede e a fronteira que o modelo estuda,
em 11 artigos; 33 descrevem uma alocação, por regra ou por estimativa, e
seis reconhecem um desvio. As restantes 89 não dizem nada sobre isso. A
origem do dado é, na maior parte, um cálculo a partir de dados de processo
(45 entradas) e só em 23 um contador dedicado. Ler o silêncio como
correspondência exata seria generoso: é precisamente o que a verificação
corrigiu em 12 entradas, onde a frase citada dizia de onde vinha o dado e
nada sobre onde estava a fronteira.

A validade temporal é a segunda. Em 89 dos 140 modelos, e em 26 artigos, o
texto não menciona que o modelo seja alguma vez refeito; em 28 não se
consegue perceber. Só 23 modelos, em oito artigos, têm um regime declarado:
por evento com gatilho explícito (16), janela móvel (3), atualização
contínua (2) ou periodicidade fixa (2). O domínio de validade não é
declarado em 121 modelos; é-o qualitativamente em oito e com limites numéricos
em 11, estes em três artigos. E a validação fora da amostra, com resultado reportado, aparece em 13
modelos de quatro artigos, contra 49 modelos com ajuste reportado nos dados
de treino. Outros seis modelos, em três artigos, descrevem uma partição em
treino e teste sem darem o resultado no teste: \textcite{Velazquez2013}
separam os registos em dois conjuntos e reportam um único R² sem dizer em
qual, \textcite{Khalid2024b} reportam só as métricas do treino e
\textcite{Tarasovskii2014} descrevem o procedimento de forma genérica. Não entram nos 13. Dos 13, 12 usam uma
partição dos dados que é aleatória~\cite{Moghadasi2025,ElMaghraby2024} ou
cuja ordem temporal não é descrita~\cite{Hamedi2019}, e o décimo
terceiro compara uma simulação com a operação real~\cite{Almanza2025}.
Nenhum descreve um teste em blocos de tempo posteriores ao ajuste. Dizer que um modelo tem R² de 0,9 nos dados que o ajustaram é o
que a literatura mais faz; testar se ainda prevê no ano seguinte, o que
menos faz.

\subsection{SQ2: o que corre mal ao modelo}
\label{sec:b-mecanismos}

Os artigos descrevem 130 ocorrências de mecanismos, em 36 dos 39. A
Tabela~\ref{tab:b-mecanismos} cruza o efeito sobre o modelo com o estatuto
da afirmação, e é esse cruzamento que muda a leitura.

{\footnotesize
\ntzebra
\begin{xltabular}{\textwidth}{@{}>{\raggedright\arraybackslash}X r r r r r@{}}
\caption[Mecanismos por efeito e estatuto da afirmação]{Mecanismos no
Corpus~B por efeito sobre o modelo e estatuto da afirmação (ocorrências;
artigos entre parênteses). \emph{General}: afirmação sobre métodos em
princípio. Aplicada: dita sobre o modelo do estudo. Demonstrada: mostrada
com os dados do estudo. A última coluna conta os artigos com pelo menos uma
ocorrência não \emph{general}.}
\label{tab:b-mecanismos}\\
\hiderowcolors
\toprule
\textbf{Efeito sobre o modelo} & \textbf{Todas} & \textbf{\emph{General}} & \textbf{Aplicada} & \textbf{Demonstr.} & \textbf{Art. sem \emph{general}} \\
\midrule
\endfirsthead
\toprule
\textbf{Efeito sobre o modelo} & \textbf{Todas} & \textbf{\emph{General}} & \textbf{Aplicada} & \textbf{Demonstr.} & \textbf{Art. sem \emph{general}} \\
\midrule
\endhead
\midrule
\rowcolor{white}\multicolumn{6}{r}{\emph{continua na página seguinte}}\\
\endfoot
\bottomrule
\endlastfoot
\showrowcolors
    Indicador sem significado tal como definido & 24 (18) & 10 & 8 & 6 & 13 \\
    Má especificação estatística & 21 (13) & 13 & 3 & 5 & 7 \\
    Comparabilidade entre instalações & 15 (11) & 6 & 7 & 2 & 8 \\
    Outro & 13 (10) & 1 & 9 & 3 & 9 \\
    Grandeza não medida & 10 (9) & 0 & 9 & 1 & 9 \\
    Consumo de base ou não produtivo & 10 (7) & 5 & 0 & 5 & 3 \\
    Dados em falta, errados ou descontínuos & 9 (6) & 1 & 5 & 3 & 5 \\
    Não estacionariedade da relação & 8 (7) & 2 & 5 & 1 & 6 \\
    Alocação entre produtos & 8 (6) & 5 & 3 & 0 & 2 \\
    Deriva ou caducidade da referência & 7 (6) & 5 & 1 & 1 & 2 \\
    Domínio de validade excedido & 4 (3) & 1 & 2 & 1 & 3 \\
    Série demasiado curta & 1 (1) & 1 & 0 & 0 & 0 \\
    \midrule
    Total & 130 (36) & 50 (19) & 52 (28) & 28 (15) & 34 \\
\end{xltabular}
}

Trinta e oito por cento das ocorrências (50 em 130) são afirmações \emph{general}:
uma definição de heteroscedasticidade, a cautela de que os processos não são
lineares, a observação de que a correlação é sensível a valores extremos. Não
são resultados dos estudos. Contá-las com as outras mede o que a literatura
repete; excluí-las mede o que os estudos encontraram, e as duas ordens não
coincidem. A má especificação estatística é o segundo efeito mais citado e
cai para o meio da tabela quando as afirmações \emph{general} saem: 13 das
suas 21 ocorrências são cautelas de manual. A deriva da referência e a
alocação entre produtos descem para dois artigos cada. No sentido oposto,
a grandeza não medida nunca é afirmação de manual: as dez ocorrências, em
nove artigos, dizem respeito ao próprio caso.\footnote{As citações dos
estudos do Corpus~B são traduções do autor.} ``Estes fatores não são todos
mensuráveis'', escrevem \textcite{Hamedi2019} sobre as variáveis que
faltavam aos seus modelos; ``o maior desafio é a disponibilidade de
instrumentação'', diz um gestor citado por \textcite{Sivill2013}; e
\textcite{EspinelBlanco2020} explicam um valor muito baixo no caso do carvão
por ``não haver medição do vapor que entra na máquina''. É o mesmo efeito
que o permutador não medido de Sines expõe no Capítulo~\ref{cha:introducao}:
uma referência não pode normalizar pelo que não se mede.

O protocolo registou a expectativa de que a grandeza não medida, a alocação
entre produtos e a deriva da referência dominassem. Não dominam. Com todas as
ocorrências, lideram o indicador sem significado tal como definido, a má
especificação e a comparabilidade entre instalações; sem as afirmações
\emph{general}, a grandeza não medida sobe para o grupo da frente, mas a
alocação e a deriva ficam no fundo. Dos três efeitos previstos, só um
se confirma, e só depois de retirar as afirmações \emph{general}.
Este resultado é uma ordem, não uma proporção: com a cobertura executada
não se pode dizer quantos dos 40 artigos reportam cada efeito, só como se
ordenam no que foi extraído.

O segundo eixo quase não devolve informação. Em 89 das 130 ocorrências
(68\%) o artigo não nomeia nenhuma causa física ou operacional, e em mais 16
a causa não cabe nas categorias. As causas nomeadas somam 25 ocorrências:
mudança de gama ou grau de produto (12, em nove artigos), qualidade da
alimentação, condições ambientais, mudança de modo de operação e alteração
da instalação (três cada) e degradação de equipamento (uma). No texto
extraído dos 39 registos não se encontrou nenhuma ocorrência de desativação
de catalisador nem de incrustação. A única frase do corpus que nomeia
incrustação~\cite{Morales2018} enumera-a entre as razões por que o
desempenho de uma caldeira se degrada com o tempo, como motivação para
monitorizar; na decisão de codificação de 18 de setembro de 2026 foi
classificada como motivação e não como limitação do modelo, e a categoria
fica, por isso, a zero. Não se afirma
que o corpus esteja livre delas. Há duas leituras, e as duas têm de ficar. A literatura
relata sintomas no modelo e raramente diz o que, no processo, os causou; ou
o instrumento, com uma lista de causas pensada a partir da refinação, não
apanha as causas que os outros setores descrevem. As 16 ocorrências de causa
``outra'' (a taxa de operação, a idade das máquinas, fatores económicos)
apoiam em parte a segunda leitura.

Por isso vale a pena ler os poucos casos com causa nomeada e resultado
mostrado. \textcite{Giacone2012} mostram que a eficiência marginal de um
forno de fusão muda com a prática de fusão sequencial.
\textcite{Beisheim2020} mostram que efeitos sazonais, como a utilização e a
temperatura ambiente, alteram o valor absoluto da intensidade energética de
um complexo integrado sem que o desempenho mude. E \textcite{Boyd2013}
mostram, ao reestimar um indicador de referência para o cimento, que a maior
parte da redução da dispersão se devia à melhoria da própria indústria, o
que tornava a referência antiga pouco representativa. No lado das afirmações
aplicadas, \textcite{Andersson2021} descrevem uma fábrica de pasta que define
vários estados normais de operação, um por qualidade de pasta, e
\textcite{Gonzalez2012} traçam a linha-alvo com o mesmo declive da referência
``porque se supõe que não há mudanças tecnológicas''. É o pressuposto que a
Secção~\ref{sec:estabilidade} testa em Sines.

\subsection{SQ3: o que se propõe}
\label{sec:b-adaptacoes}

Há 88 ocorrências de adaptações em 36 dos 39 artigos
(Tabela~\ref{tab:b-adaptacoes}). O protocolo previa que dois terços dos
artigos não propusessem nada; aconteceu o contrário, e só três não propõem
nenhuma adaptação. A expectativa estava errada, mas o número também diz
alguma coisa sobre o instrumento: 35 das 88 ocorrências caem em ``outra'',
o que mostra uma lista de tipos curta. Nessa categoria estão, por exemplo,
a agregação hierárquica de indicadores e referências de processo até ao
nível do complexo~\cite{Beisheim2020}, a carta CUSUM para garantir que o
modelo se refere ao período recente e estável~\cite{Giacone2012}, o
agrupamento prévio de empresas antes de comparar indicadores~\cite{Bonfa2017}
e a atualização periódica de um indicador de referência
setorial~\cite{Boyd2013,Boyd2017}.

\needspace{20\baselineskip}
{\footnotesize
\ntzebra
\begin{xltabular}{\textwidth}{@{}>{\raggedright\arraybackslash}X r r r r r r r r@{}}
\caption[Adaptações por tipo, estado e enquadramento]{Adaptações no
Corpus~B por tipo, estado e enquadramento normativo (ocorrências; artigos
entre parênteses). Ligadas: o artigo liga explicitamente a adaptação a um
mecanismo que descreveu. \emph{Implícito} exige que a adaptação corresponda
a um conceito da norma --- variável relevante, fator estático, revisão da
referência, fronteira ---, e não que o parágrafo use vocabulário da norma.}
\label{tab:b-adaptacoes}\\
\hiderowcolors
\toprule
& & \multicolumn{3}{c}{\textbf{Estado}} & \multicolumn{3}{c}{\textbf{Enquadramento}} & \\
\cmidrule(lr){3-5}\cmidrule(lr){6-8}
\textbf{Tipo} & \textbf{Total} & \textbf{Prop.} & \textbf{Impl.} & \textbf{Aval.} & \textbf{ISO expl.} & \textbf{Implíc.} & \textbf{Ausente} & \textbf{Ligadas} \\
\midrule
\endfirsthead
\toprule
& & \multicolumn{3}{c}{\textbf{Estado}} & \multicolumn{3}{c}{\textbf{Enquadramento}} & \\
\cmidrule(lr){3-5}\cmidrule(lr){6-8}
\textbf{Tipo} & \textbf{Total} & \textbf{Prop.} & \textbf{Impl.} & \textbf{Aval.} & \textbf{ISO expl.} & \textbf{Implíc.} & \textbf{Ausente} & \textbf{Ligadas} \\
\midrule
\endhead
%% sem rodape de continuacao: o \needspace acima garante que nao parte
\bottomrule
\endlastfoot
\showrowcolors
    Outra & 35 (23) & 20 & 11 & 4 & 3 & 13 & 19 & 16 \\
    Variável de normalização adicional & 17 (15) & 4 & 11 & 2 & 4 & 8 & 5 & 13 \\
    Exclusão de valores atípicos & 10 (10) & 0 & 10 & 0 & 0 & 3 & 7 & 4 \\
    Referência segmentada & 9 (8) & 3 & 6 & 0 & 0 & 6 & 3 & 2 \\
    Medição melhorada & 7 (6) & 6 & 1 & 0 & 0 & 2 & 5 & 4 \\
    Normalização a produto de referência & 5 (5) & 1 & 4 & 0 & 0 & 5 & 0 & 2 \\
    Reestimação contínua & 2 (2) & 1 & 1 & 0 & 0 & 2 & 0 & 2 \\
    Estimação ponderada & 2 (2) & 0 & 2 & 0 & 0 & 0 & 2 & 0 \\
    Janela móvel & 1 (1) & 0 & 1 & 0 & 0 & 1 & 0 & 0 \\
    \midrule
    Total & 88 (36) & 35 & 47 & 6 & 7 & 40 & 41 & 43 \\
\end{xltabular}
}

Propor e testar são coisas diferentes, e a distância entre as duas é o
resultado mais útil desta secção. Das 88 adaptações, 35 são propostas para
trabalho futuro e 47 foram implementadas no próprio estudo; só seis foram
comparadas com uma alternativa. \textcite{Velazquez2013} são dos poucos que
o fazem numa refinaria: ajustar a referência de uma unidade de reformação às
variáveis de influência dá uma poupança de \revisto{2,82~GWh, contra 2,03~GWh com uma referência constante; este último valor é 28\% inferior ao primeiro.} Os outros cinco comparam uma
rede neuronal com a regressão~\cite{Hamedi2019}, a fronteira estocástica com
os mínimos quadrados~\cite{Boyd2017}, a referência atualizada com a
original~\cite{Boyd2013}, um indicador com preços ponderados com o
simples~\cite{Scripcariu2021} e um novo indicador que junta informação do
processo com o anterior~\cite{Khalid2024b}. Só duas das seis usam dados que
não serviram para ajustar o modelo: a comparação de \textcite{Hamedi2019},
feita num conjunto de teste, e a de \textcite{Boyd2013}, que aplica o modelo
antigo aos dados novos.

As adaptações que tocam a primeira lacuna quase não saem do papel: melhorar
a medição é proposto seis vezes e feito uma. As que tocam a validade
temporal são raras: dois casos de reestimação contínua e um de janela
móvel, \revisto{contra dez exclusões de valores atípicos, uma das adaptações implementadas mais frequentes, embora retirar observações possa alterar o ajuste sem alterar a sua forma funcional.}

O enquadramento normativo explícito é raro, como o protocolo previa: sete
das 88 adaptações, em cinco artigos, justificam-se citando a ISO~50001 ou a
ISO~50006. Outras 40 são implícitas e 41 não têm enquadramento nenhum. Esta
repartição depende inteiramente do critério de ``implícito'', e o adotado é o
estrito: conta a adaptação que corresponde a um conceito da norma --- uma
variável relevante, um fator estático, a revisão da referência, a fronteira
do indicador --- sem a citar. O critério largo, o de contar a presença de
vocabulário da norma no parágrafo, classificaria quase todas como implícitas:
a palavra EnPI está em quase todo o corpus porque a expressão de pesquisa~B a
exige, e o campo deixaria de medir alguma coisa. Outra expectativa registada, a de que as adaptações raramente se
ligam ao mecanismo que tratam, tem de ser reformulada: 43 das 88 ligam-se
explicitamente a um mecanismo descrito no mesmo artigo. ``Raramente'' é
errado; metade não se liga, e é isso que se pode dizer.

\subsection{Limitações, barreiras e reprodutibilidade}

Vinte e quatro artigos declaram 46 limitações do próprio estudo, sobretudo
acesso a dados (16), âmbito (14) e cobertura da amostra (13). Vinte e oito
descrevem 62 barreiras à implementação. Estes dois campos são texto livre,
não passaram pela verificação e a leitura que se segue é interpretativa. As
barreiras organizacionais dominam: compromisso da gestão, tempo e formação
das equipas, prioridade da energia face à produção. Mas há um segundo grupo,
mais próximo desta dissertação, em que a barreira é o próprio método ou a
falta de medição. \textcite{Hamedi2019} notam que as empresas continuam a
usar o consumo específico simples ``pela simplicidade e pelas exigências das
normas nacionais e das auditorias''; \textcite{Beisheim2020} duvidam que uma
empresa aceite o esforço de modelação ``só para calcular referências
energéticas''; \textcite{Bruni2021} descrevem a fase de normalização como
demorada e dependente de pessoal muito especializado; \textcite{Bonfa2017}
encontram a maioria das empresas inquiridas sem sistema de medição do ar
comprimido; e \textcite{Boyd2013} referem dados de instalações tratados como
confidenciais.

A reprodutibilidade fica a meio. As equações do modelo estão no artigo em 24
dos 39 casos e em parte em mais nove; os coeficientes, em 12 e em parte em
13. Dados utilizáveis estão disponíveis em três artigos e em parte noutros
três. Nenhum dos 39 registos declara código disponível, e um declara-o em
parte. Como as ausências não foram verificadas, este zero lê-se como
``nenhum declara'', não como ``nenhum tem''.

\subsection{Os estudos mais próximos do caso}
\label{sec:corpus-b-proximos}

A SQ4 foi retirada e esta subsecção não a substitui: não avalia se os
métodos se transferem para Sines. Lê criticamente os 13 artigos com uma
aplicação empírica em refinaria ou petroquímica, porque são os vizinhos
diretos do caso e é contra eles que o contributo da dissertação se mede
(Tabela~\ref{tab:b-proximos}).

{\scriptsize
\ntzebra
\begin{xltabular}{\textwidth}{@{}>{\raggedright\arraybackslash}X>{\raggedright\arraybackslash}p{3.3cm}>{\raggedright\arraybackslash}p{1.15cm}>{\raggedright\arraybackslash}p{1.75cm}>{\raggedright\arraybackslash}p{1.55cm}>{\raggedright\arraybackslash}p{1.3cm}@{}}
\caption[Estudos do Corpus B com aplicação em refinaria ou petroquímica]{Os
13 estudos do Corpus~B com aplicação empírica em refinaria ou petroquímica.
Fronteira (C5), reestimação (C6), domínio de validade (C7) e validação fora
da amostra (C8), no modelo mais forte de cada estudo. RL: regressão linear;
n.r.: não reportado; n.a.: não aplicável.}
\label{tab:b-proximos}\\
\hiderowcolors
\toprule
\textbf{Estudo e objeto} & \textbf{Modelo e variáveis além da carga} & \textbf{Fronteira} & \textbf{Reestimação} & \textbf{Domínio} & \textbf{Fora da amostra} \\
\midrule
\endfirsthead
\toprule
\textbf{Estudo e objeto} & \textbf{Modelo e variáveis além da carga} & \textbf{Fronteira} & \textbf{Reestimação} & \textbf{Domínio} & \textbf{Fora da amostra} \\
\midrule
\endhead
\midrule
\rowcolor{white}\multicolumn{6}{r}{\emph{continua na página seguinte}}\\
\endfoot
\bottomrule
\endlastfoot
\showrowcolors
    \textcite{Velazquez2013}: unidade de reformação catalítica & RL horária; ambiente, alimentação, modo & n.r. & não mencionada & não declarado & n.r. \\
    \textcite{Yao2016}: compressão de gás de \emph{cracking} & RL; modo & n.r. & por evento & qualitativo & n.r. \\
    \textcite{Moghadasi2025}: descarbonatação de etano, refinaria de gás & Árvores de decisão (GBM), dados do DCS; modo & n.r. & não mencionada & não declarado & reportada \\
    \textcite{Beisheim2019}: processos de um complexo químico & Modelo substituto não linear; ambiente & exata & por evento & qualitativo & n.r. \\
    \textcite{Hamedi2019}: complexo petroquímico & RL e rede neuronal; ambiente, modo, produtos & n.r. & não declarada & não declarado & reportada \\
    \textcite{Beisheim2020}: complexo integrado (dados modificados) & Rácios e regressões agregados; ambiente, utilização, produtos & exata em parte & por evento & \revisto{n.r.} & n.r. \\
    \textcite{Ochoa2017}: instalação petroquímica & RL de produção equivalente; só carga & n.r. & não mencionada & não declarado & n.r. \\
    \textcite{Morales2018}: caldeira de vapor & Índice, RL, carta de controlo & n.r. & não mencionada & não declarado & n.r. \\
    \textcite{Almanza2025}: central de cogeração de refinaria & Simulação termodinâmica & exata & não mencionada & não declarado & reportada \\
    \textcite{Ikeyama2017}: instalação da Saudi Aramco (caso de fornecedor) & Modelos ajustados a modos e crude & n.r. & \revisto{n.r.} & não declarado & n.r. \\
    \textcite{Hashim2019}: empresa de petróleo e gás, várias divisões & Índice de implementação do SGE & n.r. & não declarada & não declarado & n.a. \\
    \textcite{Maggiore2020}: auditorias em vários setores & Rácio & n.r. & não mencionada & não declarado & n.a. \\
    \textcite{Hashim2022}: empresa petrolífera & Rácio; só carga & n.r. & não mencionada & não declarado & n.a. \\
\end{xltabular}
}

Os 13 dividem-se em três grupos. O primeiro modela uma unidade de processo
com dados do historiador ou do sistema de controlo: \textcite{Velazquez2013},
\textcite{Yao2016}, \textcite{Moghadasi2025} e \textcite{Beisheim2019}. São os
únicos com séries finas e os únicos que se aproximam do problema de Sines.
Nenhum junta as quatro condições da tabela. \textcite{Beisheim2019} é o que
mais se aproxima, com a fronteira descrita, uma regra para refazer o modelo
depois de alterações na instalação e um domínio de validade dado pelos
pontos usados na identificação, mas sem teste fora da amostra.
\textcite{Moghadasi2025} reportam um R² de 0,99 num conjunto de validação
com 20\% dos dados escolhidos ao acaso; numa série sub-horária de um ano,
pontos vizinhos ficam dos dois lados da partição e a validação mede
interpolação, não previsão.
\textcite{Yao2016} propõem refazer o modelo quando o consumo cai três
desvios-padrão abaixo da referência. A regra tem um problema que atravessa
esta dissertação: uma melhoria real do desempenho dispara a mesma condição
que um modelo caducado, e refazer a referência nesse momento apaga a
poupança que se queria medir.

O segundo grupo não modela o desempenho de uma unidade.
\textcite{Almanza2025} simulam uma central de cogeração em três cenários, e
o modelo subestima o consumo de gás natural face aos valores reais em todas
as caldeiras menos uma. \textcite{Hashim2019} constroem um indicador do
estado de implementação do sistema de gestão, que os próprios autores dizem
não medir o desempenho energético real. \textcite{Ikeyama2017} descrevem
normas e citam um caso de fornecedor sem dados. \textcite{Maggiore2020} e
\textcite{Hashim2022} usam rácios em auditorias e programas de gestão.
Estão no Corpus~B com toda a legitimidade, mas não informam sobre
referências de processo.

O terceiro grupo modela instalações petroquímicas inteiras ou equipamentos:
\textcite{Hamedi2019}, \textcite{Ochoa2017}, \textcite{Morales2018} e
\textcite{Beisheim2020}. \textcite{Hamedi2019} é o único que compara duas
famílias de modelos com validação fora da amostra, e o único que reconhece,
sobre o próprio caso, que parte dos fatores que movem o consumo não é
mensurável. \textcite{Ochoa2017} normalizam só pela produção equivalente. Dos quatro,
só \textcite{Beisheim2020} dizem em que condições a referência seria
refeita.

Lidos em conjunto, os estudos mais próximos têm as mesmas lacunas que o
corpus e nenhum resolve a combinação que um gestor de energia de uma
refinaria precisa: a energia medida onde o modelo diz, a carga e as outras
variáveis relevantes medidas, um domínio de validade explícito, um teste
fora da amostra em blocos de tempo e uma regra para refazer a referência que
não confunda melhoria com caducidade.

\subsection{Alcance e limites desta \ajuste{leitura}}

Os resultados desta secção descrevem o que 39 artigos dizem, lidos por um
modelo de linguagem e relidos por outro da mesma família, com os casos
duvidosos decididos pelo autor. Não são propriedades verificadas da
literatura. As ausências não foram confirmadas por leitura das secções
relevantes, \revisto{pelo que as ocorrências identificadas podem subcontar o que os artigos contêm, enquanto as classificações ``não reportado'' podem sobrestimar as ausências.} O rastreio de referências não foi feito, e a
literatura de \emph{benchmarking} que os estudos citam ficou fora
(Secção~\ref{sec:pesquisa}). Faltam 137 textos integrais não recuperados, com
perda desigual. As listas de tipos são curtas: 35 das 88 adaptações, 16 das
130 causas e 13 dos 130 efeitos caíram em ``outra''. O Corpus~B não é uma
amostra aleatória de nada: são os estudos destes setores que estas expressões
encontraram, com cimento e siderurgia a dar quase metade das aplicações. As
ordens de efeitos e as contagens servem para comparar dentro do corpus, não
para estimar frequências na literatura.

\section{Síntese: o que falta à literatura}
\label{sec:gap}

A pergunta com que o capítulo abriu era se a literatura dá a um gestor de
energia um método para decidir se a referência de uma unidade ainda serve.
A resposta, nos limites desta revisão, é que dá peças e não o método. Os
rácios e as regressões sobre a carga estão bem descritos, a norma que os
enquadra é citada e há casos em quase todos os setores de processo contínuo.
\ajuste{Citar a norma não traz, porém, um procedimento de verificação: no
Corpus~A, 173 das 187 publicações que invocam a família ISO~50001 e têm o
campo resolvido não invocam protocolo algum de medição e verificação
(Secção~\ref{sec:sq1-sintese}). O vocabulário da norma está adotado; o quadro
para confirmar o que se mede, não.} O que falta está nas condições em que
esses modelos se podem usar como referência, e concentra-se em três lacunas
documentadas no Corpus~B (Tabela~\ref{tab:gap}).

A primeira é a medição. A energia que o modelo explica raramente é descrita
como medida na fronteira que o modelo estuda. \ajuste{Os 39 estudos dizem 163
vezes de onde vem a energia que os seus modelos usam, uma vez por vetor
energético; em 35 dessas vezes o que se mede coincide com a fronteira que o
modelo estuda, e em 89 o artigo não diz nada sobre isso.} A grandeza não medida é o
único efeito que os estudos descrevem sempre sobre o próprio caso, e \revisto{melhorar a medição é proposto seis vezes e implementado uma vez.} Um modelo cuja
variável dependente é um cálculo, numa fronteira que não coincide com a do
contador, pode ajustar bem e medir outra coisa.

A segunda é a validade no tempo. \revisto{Só 23 de 140 modelos têm um regime de reestimação declarado, 19 um domínio de validade e 13 uma validação fora da amostra;} das 88 adaptações, seis foram avaliadas contra uma alternativa e
duas com dados que não serviram para ajustar. A literatura constrói
referências e raramente as põe à prova depois. Onde propõe uma regra para as
refazer, a regra pode confundir uma melhoria real com um modelo caducado.

A terceira é a normalização. Em 78 dos 140 modelos a carga é a única
variável que normaliza o consumo, e a qualidade da alimentação aparece em
quatro artigos. \ajuste{Numa unidade de processo contínuo a simplificação tem
consequência direta: mudar de crude, de severidade ou de gama de produtos
altera a energia necessária sem alterar a carga, e um modelo que só conhece a
carga lê essa mudança como perda de desempenho. É o mesmo mecanismo que o
Capítulo~\ref{cha:diagnostico} tem de separar em Sines.} O efeito mais citado
nos estudos é, precisamente, um indicador que não significa o que devia tal
como está definido. A norma pede
que se identifiquem as variáveis relevantes; os estudos quase nunca mostram
que a carga chega.

Há ainda uma falha de outra natureza. \ajuste{Nenhum dos 39 estudos declara
código disponível e três disponibilizam os dados, pelo que nenhuma das
referências publicadas se reaplica a outra instalação sem ser reconstruída a
partir do texto: reescrever as equações, reestimar os coeficientes e repetir
escolhas de fronteira que o artigo não declara. Ao contrário das três
anteriores, esta não é uma lacuna dos modelos de desempenho energético, é uma
condição de acesso.}


\ajuste{A coluna da direita da Tabela~\ref{tab:gap} não é uma justificação
a posteriori. A
plataforma, o modelo de dados e o protocolo de avaliação foram decididos e
congelados antes de o Corpus~B ser extraído, com as datas registadas nos
repositórios (Apêndice~\ref{app:protocolo-sr}). O que a revisão mostra não é
que essas decisões estejam certas, mas que os problemas encontrados em Sines
não são uma particularidade do caso: são os mesmos que a literatura descreve
e não resolve.}

O Capítulo~\ref{cha:introducao} resumia o problema numa frase: mede-se tudo e
sabe-se tarde. A literatura revista repete o padrão. Mede, com rácios e
regressões sobre a carga e coeficientes de determinação calculados nos dados
de ajuste. Saber se a referência ainda vale, onde foi medida, até onde se
estende e se prevê o ano seguinte, é o que quase nenhum dos 39 estudos mostra.
É esse o espaço que os capítulos seguintes ocupam.


\needspace{24\baselineskip}
{\footnotesize
\ntzebra
\begin{xltabular}{\textwidth}{@{}>{\raggedright\arraybackslash}p{2.6cm}>{\raggedright\arraybackslash}X>{\raggedright\arraybackslash}p{4.6cm}@{}}
\caption[Lacunas documentadas e resposta da dissertação]{Lacunas documentadas
no Corpus~B e onde a dissertação lhes responde. \ajuste{As decisões dos
Capítulos~\ref{cha:metodos} a~\ref{cha:diagnostico} antecedem a extração do
Corpus~B, de 12 a 15 de setembro de 2026. A última linha é de outra natureza:
não é uma lacuna dos modelos, é uma condição de acesso à literatura.}}
\label{tab:gap}\\
\hiderowcolors
\toprule
\textbf{Lacuna} & \textbf{Evidência no Corpus~B} & \textbf{Resposta nesta dissertação} \\
\midrule
\endfirsthead
\toprule
\textbf{Lacuna} & \textbf{Evidência no Corpus~B} & \textbf{Resposta nesta dissertação} \\
\midrule
\endhead
%% sem rodape de continuacao: o \needspace acima garante que nao parte
\bottomrule
\endlastfoot
\showrowcolors
    Energia medida na fronteira do modelo &
    \ajuste{35 de 163 descrições de medição com correspondência exata, 89 sem
    informação; grandeza não medida descrita sempre sobre o próprio caso
    (9 artigos); medição melhorada proposta 6 vezes, feita 1} &
    Reconciliação das fontes e tabela de mapeamento como contrato
    (Secções~\ref{sec:fontes} e~\ref{sec:modelo-dimensional}); cobertura por
    unidade e vetor (Secção~\ref{sec:cobertura}); perfil dos dados
    (Secção~\ref{sec:perfil-dados}) \\
    \addlinespace
    Validade da referência no tempo &
    \revisto{Regime de reestimação em 23 de 140 modelos (8 artigos); domínio de validade em 19 (numérico em 11, 3 artigos);} validação fora da amostra em
    13 (4 artigos); 6 de 88 adaptações avaliadas &
    Protocolo pré-registado com replicação e comparadores triviais
    (Secção~\ref{sec:protocolo-baselines}); desempenho preditivo, estabilidade,
    calibração e detetabilidade
    (Secções~\ref{sec:poder-preditivo} a~\ref{sec:detetabilidade}) \\
    \addlinespace
    Normalização só pela carga &
    Carga como única variável em 78 de 140 modelos; qualidade da alimentação
    em 4 artigos; \ajuste{indicador sem significado tal como definido é o efeito
    mais citado (24 casos, e 13 artigos descrevem-no sobre o próprio caso)} &
    Métricas e referência em produção (Secção~\ref{sec:metricas});
    suficiência da variável explicativa (Secção~\ref{sec:suficiencia}) \\
    \addlinespace
    \ajuste{Reprodutibilidade (transversal)} &
    \ajuste{Código declarado disponível em 0 de 39 estudos; dados em 3} &
    \ajuste{Manifesto e proveniência por corrida
    (Secção~\ref{sec:reprodutibilidade}); pacote público com dados sintéticos
    como trabalho futuro} \\
\end{xltabular}
}